\section{Cross-lingual Self-supervised Speech Representation Learning}
\label{sec:pretrain}

As a first step, we train a self-supervised model of speech representations on the data outlined above as well as other existing public corpora. 
We use wav2vec 2.0 for pre-training on unlabeled data which we later use as the basis for several downstream speech tasks (\citealt{baevski2020wav};~\textsection\ref{sec:pretrain-w2v}).
The resulting models were pre-trained on \numptlang{} languages which is over four times the number of languages of known prior work~(\citealt{zhang2023usm};~\textsection\ref{sec:pretrain-setup}).
The increased language coverage results in better performance for both ASR and LID compared to XLS-R (\citealt{babu22_interspeech};~\textsection\ref{sec:pretrain-results}) which covered 128 languages and is publicly available.


\subsection{Method: wav2vec 2.0 and XLS-R}
\label{sec:pretrain-w2v}

Our work builds on~\citet{babu22_interspeech} who pretrain wav2vec models on data from multiple languages.
The wav2vec project created a series of models for learning self-supervised speech representations~\citep{schneider2019wav2vec,baevski2019vqwav2vec,baevski2020wav} from unlabeled speech data.
The resulting models can then be used to solve downstream speech tasks by fine-tuning them on labeled data or by tackling these tasks without labeled data using unsupervised learning~\citep{baevski2021unsupervised,liu2022endtoend}.

The most prominent one is wav2vec 2.0~\citep{baevski2020wav} which enables building speech recognition models with only ten minutes of labeled data and even no labeled data at all~\citep{baevski2021unsupervised}.
The basic architecture of wav2vec 2.0 is as follows: a convolutional feature encoder $f: \Inp \mapsto \Feat$ maps raw audio~$\Inp$ to latent speech representations $\ze_1, \dots, \ze_T$ which are input to a Transformer $g: \Feat \mapsto \Context$ to output context representations $\cc_1, \dots, \cc_T$~\citep{baevski2019vqwav2vec}.
Each $\ze_t$ represents 25ms of audio strided by 20ms and the Transformer architecture follows BERT~\citep{vaswani2017transformer,devlin2018bert}.

During training the feature encoder representations are discretized to $\zq_1, \dots, \zq_T$ with a quantization module $\Feat \mapsto \QFeat$ to represent the targets in the objective.
The quantization module uses a Gumbel softmax to choose entries form the codebooks and the chosen entries are concatenated~\citep{jegou2011ieee,jang2016gumbel,baevski2019vqwav2vec}.

The model is trained by solving a contrastive task over masked feature encoder outputs.
At training time, spans of ten time steps with random starting indices are masked.
The objective requires identifying the true quantized latent $\zq_t$ for a masked time-step within a set of $K=100$ distractors sampled from other masked time steps.
The objective is augmented by a codebook diversity penalty to encourage the model to use all codebook entries~\citep{dieleman2018challenge}.

XLSR and XLS-R train wav2vec 2.0 on many different languages from several datasets to obtain cross-lingual representations~\citep{conneau2020unsupervised,babu22_interspeech}.
In order to balance the training data, two data sampling steps are performed.
First, for each dataset, we sample the data for the different languages $L$ from a distribution $p_l\sim\left(\frac{n_l}{N}\right)^{\beta_L}$ where $l=1,\ldots,L$, $n_l$ is the amount of unlabeled data for each language in the dataset, $N$ is the total amount of training in the dataset, and $\beta_L$ is the upsampling factor which controls the trade-off between high- and low-resource languages during pretraining.
Second, we balance the different datasets by treating each dataset as a language in the above sampling scheme with a sampling parameter $\beta_D$.


\insertPretrainSetup

\subsection{Pre-training Setup}
\label{sec:pretrain-setup}

\paragraph{Hyperparameters.}
We largely follow prior work in training cross-lingual wav2vec 2.0 models~\citep{conneau2020unsupervised,babu22_interspeech} and use the wav2vec 2.0 implementation available in fairseq~\citep{ott2019fairseq} to train models with roughly 300M and 1B parameters (\autoref{tab:pretrain-models}).
To make efficient use of GPU memory, we use a fully sharded backend~\citep{rajbhandari2021zero} as well as activation checkpointing~\citep{chen20216checkpointing} implemented in FairScale~\citep{FairScale2021}. 

Our models are optimized with Adam~\citep{kingma2015adam} and the learning rate is warmed up for the first 32K steps followed by polynomial decay to zero for the remainder of training. 
Training audio sequences are cropped to a maximum of 320K samples, or 20 seconds, and all models were pre-trained for a total of one million updates on A100 GPUs with 80GB of memory.
The \modelnameb{0.3} model was trained with an effective batch size of 2.3 hours of data across 48 GPUs and the \modelnameb{1} model was trained with an effective batch size of 3.5 hours on 64 GPUs.

\paragraph{Data.}
The pre-training data covers about 491K hours in \numptlang{} languages. 
This data is drawn from six training corpora with different characteristics, including the corpora used in XLS-R~\citep{babu22_interspeech}: 


\begin{itemize}
    \item \MMSU: \nummmsulang{} languages comprising \nummmsuhrs{} hours (\textsection\ref{sec:data-paired}). 
    \item Multilingual Librispech (MLS): 8 European languages of read books totaling 50K hours~\citep{pratap2020mls}
    \item CommonVoice (CV): 89 languages totaling 8.8 hours of read Wikipedia text; we use v9.0 of the corpus~\citep{ardila2019common})
    \item VoxLingua-107 (VL): 107 languages totaling 5.3K hours of YouTube content~\citep{valk2020voxlingua}
    \item BABEL (BBL): 17 African and Asian languages totaling about 1K hours of conversational telephone data~\citep{gales2014babel}
    \item VoxPopuli (VP): 371K hours of unlabeled speech data in 23 languages derived from European Parliament event recordings~\citep{wang2021voxpopuli}

\end{itemize}

Pre-training only uses the speech audio and none of the transcriptions and we balance the data following the strategy outlined in~\textsection\ref{sec:pretrain-w2v} using $\beta_L = \beta_D = 0.5$


\subsection{Comparison to XLS-R}
\label{sec:pretrain-results}

\insertPretrainComparisonASR

To better understand how the {\modelname} models compare to XLS-R we fine-tuned both for automatic speech recognition on the 61 languages of the FLEURS benchmark for which \MMS{} provides training data.
Models are evaluated without a language model and we report the average character error rate (CER) on FLEURS development data over all languages.

\autoref{fig:pretrain-comparison-asr} shows that the {\modelname} models perform better: for the 300B size, {\modelname} has 0.6 lower CER than XLS-R, and for the 1B size, the difference is 0.7 CER.
More capacity helps to improve performance: 
for XLS-R the error rate decreases by 3.2 CER absolute when scaling the number of parameters from 300M to 1B and for \MMS{} there is a 3.0 CER improvement.

\modelname{} pre-trains on over ten times the number of languages of XLS-R and this improves performance, particularly on low resource languages (\autoref{fig:pretrain-xlsr-mms-perlang}) such as Amharic (amh), Lao (lao) or Malayalam (mal).
Compared to XLS-R, the pre-training data of \modelname{} covers the following languages which improve: Chewa (nya), Fulah (ful) and Oromo (orm).
However, improvements at low resource languages result in a small degradation in some of the high-resource languages such as English (eng) or Spanish (spa) but there are also other languages such as Tajik (tgjk) or Welsh (cym) which perform less well.

\insertPretainXLSRvsMMS

Equipped with this new pre-trained model we now investigate the use of \MMS{} and \GRN{} for several downstream tasks.
